% Part: intuitionistic-logic
% Chapter: tableaux
% Section: introduction

\documentclass[../../../include/open-logic-section]{subfiles}

\begin{document}

\olfileid{int}{tab}{int}

\olsection{పరిచయం}

\tetoken{టాబ్లోలు}{!!^{tableau}s} అనేవి
\tetoken{చిహ్నిత
  సూత్రాలతో}{!!{signed
  formula}s} ఏర్పడిన కొన్ని (కిందికి శాఖలుగా
విడిపోయే) వృక్షాలు. చిహ్నిత సూత్రం అంటే
సత్యమూల్య చిహ్నం ($\True$ లేదా
$\False$), \tetoken{ఒక వాక్యం}{!!a{sentence}}
కలిసిన జత:
\[
\sFmla{\True}{!A} \text{ లేదా } \sFmla{\False}{!A}.
\]
\tetoken{ఒక టాబ్లో}{!!^a{tableau}}
కొన్ని \emph{పరికల్పనలతో} మొదలవుతుంది.
తరువాతి ప్రతి \tetoken{చిహ్నిత
సూత్రం}{!!{signed formula}}
నిగమన నియమాల్లో ఒకదాన్ని వర్తింపజేయడం
వల్ల ఏర్పడుతుంది. కొన్ని నియమాలు వృక్షపు
చివరి నోడ్‌కు ఒకటి లేదా అంతకంటే ఎక్కువ
\tetoken{చిహ్నిత సూత్రాలను}{!!{signed formula}s}
చేరుస్తాయి; మరికొన్ని రెండు కొత్త చివరి
నోడ్‌లను చేర్చి రెండు శాఖలను ఏర్పరుస్తాయి.
ఒక నియమాన్ని వర్తింపజేసిన
\tetoken{చిహ్నిత సూత్రం}{!!{signed formula}}
కంటే ఫలిత \tetoken{చిహ్నిత
సూత్రాల్లోని}{!!{signed formula}s}
\tetoken{సూత్రం}{!!{formula}} తక్కువ సంక్లిష్టం.
ఒక శాఖలో $\sFmla{\True}{!A}$,
$\sFmla{\False}{!A}$ రెండూ ఉంటే,
ఆ శాఖను \emph{సంవృతం} అంటాం.
\tetoken{ఒక టాబ్లోలోని}{!!a{tableau}}
ప్రతి శాఖ సంవృతమైతే, మొత్తం
\tetoken{టాబ్లో}{!!{tableau}} సంవృతం.
సంవృత \tetoken{టాబ్లో}{!!{tableau}},
\tetoken{టాబ్లోను}{!!{tableau}}
ప్రారంభించడానికి వాడిన
\tetoken{చిహ్నిత సూత్రాల}{!!{signed formula}s}
సమితి అసంతృప్తిపరచదగినదని చూపే
\tetoken{ఒక వ్యుత్పత్తి}{!!a{derivation}}.
దీని ద్వారా $\Proves$ సంబంధాన్ని
నిర్వచించవచ్చు: $\Gamma \Proves !A$
అయితే, అప్పుడే $\Gamma_0 = \{!B_1,
\dots, !B_n\} \subseteq \Gamma$
అయ్యే ఏదో పరిమిత సమితికి,
కింది పరికల్పనలపై సంవృత
\tetoken{టాబ్లో}{!!{tableau}} ఉంటుంది:
\[
\{\sFmla{\False}{!A}, \sFmla{\True}{!B_1}, \dots, \sFmla{\True}{!B_n}\}.
\]

అంతఃప్రజ్ఞావాద తర్కం కోసం
\tetoken{చిహ్నిత సూత్రం}{!!{signed formula}}
భావనను విస్తరించి, సంధాయకాల నియమాలను
సవరించాలి. చిహ్నం ($\True$ లేదా
$\False$)తో పాటు, మోడల్
\tetoken{టాబ్లోల్లోని}{!!{tableau}s}
\tetoken{సూత్రాలకు}{!!{formula}s}
\emph{పూర్వసూచికలు}~$\sigma$ కూడా ఉంటాయి.
ఆ పూర్వసూచికలు ధన పూర్ణసంఖ్యల
శూన్యం కాని క్రమాలు; అంటే
$\sigma \in
(\PosInt)^* \setminus \{\emptyseq\}$.
వాటిని వ్రాసేటప్పుడు చుట్టూ ఉండే
$\tuple{\ }$ను వదిలి, విడివిడిగా ఉన్న
\tetoken{మూలకాలను}{!!{element}s}
$,$లకు బదులు~$.$లతో వేరు చేస్తాం.
$\sigma$ ఒక పూర్వసూచిక అయితే
$\sigma.n$ అనేది $\sigma \concat
\tuple{n}$; ఉదాహరణకు $\sigma = 1.2.1$
అయితే $\sigma.3$ అనేది $1.2.1.3$.
అందువల్ల, ఉదాహరణకు,
\[
\sFmla{\True}{!A \lif (!B \lif !C)}[1.2]
\]
అనేది \emph{పూర్వసూచిక గల
\tetoken{చిహ్నిత సూత్రం}{!!{signed formula}}}
(లేదా సంక్షిప్తంగా \emph{పూర్వసూచిక గల
  \tetoken{సూత్రం}{!!{formula}}}).

అంతఃప్రజ్ఞాత్మకంగా, టాబ్లో శాఖపై
ఉన్న \tetoken{సూత్రాలను}{!!{formula}s}
సంతృప్తిపరచవచ్చే నమూనాలోని ఒక
లోకానికి పూర్వసూచిక పేరు. ఒక
\tetoken{టాబ్లో}{!!a{tableau}}లో
$\sigma$ ఏదో లోకానికి పేరైతే,
$\sigma.n$ అనేది~$\sigma$ పేరుగల
లోకం నుంచి ప్రాప్యమైన లోకానికి పేరు.

అంతఃప్రజ్ఞావాద నమూనాల్లో ప్రాప్యత
సంబంధం స్వావర్తన, సంక్రామక.
పూర్వసూచికల పరంగా, $\sigma$
అనేది~$\sigma$ నుంచే ప్రాప్యం;~$\sigma$ను
విస్తరించే ఏ పూర్వసూచికైనా, అంటే
$\sigma.n_1.\cdots.n_k$ రూపంలోనిదైనా,
దాని నుంచి ప్రాప్యమే. $\sigma$ను,
దాని ప్రతి విస్తరణనూ సూచించడానికి
$\sigma.*$ అనే సంకేతాన్ని వాడుదాం.
మరో మాటలో, $\sigma.*$ పూర్వసూచికలే
~$\sigma$ నుంచి ప్రాప్యమైన అన్ని,
ఆ పూర్వసూచికలు మాత్రమే.

\end{document}
